\documentclass[a4paper]{article}

% ---------- Encoding, lingua, font ----------
\usepackage[utf8]{inputenc}
\usepackage[T1]{fontenc}
\usepackage{textcomp}
\usepackage[english]{babel}
\renewcommand{\familydefault}{\sfdefault}
\emergencystretch=1.5em
\usepackage{enumitem}   % nel preambolo

% ---------- Colori (UNA sola volta, prima di tikz/tcolorbox) ----------
\usepackage[dvipsnames,table]{xcolor}
\definecolor{head}{HTML}{DCE6F2}
\definecolor{rowgray}{HTML}{F2F2F2}

% ---------- Matematica ----------
\usepackage{amsmath, amssymb, amsfonts, amsthm}
\usepackage{mathtools}
\DeclarePairedDelimiter\ceil{\lceil}{\rceil}
\DeclarePairedDelimiter\floor{\lfloor}{\rfloor}
\usepackage{cancel}
\usepackage{gensymb}

% ---------- Tabelle, figure, utilità ----------
\usepackage{booktabs, xltabular, tabularx, makecell}
\usepackage{subcaption}
\usepackage{float}
\usepackage{graphicx}
\usepackage{import}
\usepackage{pdfpages}
\usepackage{xifthen}
\pdfminorversion=7
\pdfsuppresswarningpagegroup=1
\graphicspath{{./figures/}}

\renewcommand{\arraystretch}{1.5}
\newcolumntype{Y}{>{\raggedright\arraybackslash}X}

% ---------- TikZ / pgfplots ----------
\usepackage{tikz}
\usetikzlibrary{shapes, arrows, automata, petri, decorations.markings,
                decorations.pathreplacing, positioning, calc, quotes}
\usepackage{circuitikz}
\usepackage{forest}
\usepackage{pgfplots}
\pgfplotsset{compat=1.18, width=10cm}

\tikzset{
    place/.style={circle, thick, draw=black, minimum size=6mm},
    state/.style={circle, thick, draw=black, fill=white, minimum size=6mm},
}

% ---------- Palette (solo quella nuova) ----------
\definecolor{thmcol}{HTML}{2563EB}   % blu     - teoremi
\definecolor{lemcol}{HTML}{7C3AED}   % viola   - lemmi
\definecolor{defcol}{HTML}{059669}   % verde   - definizioni
\definecolor{excol}{HTML}{EA580C}    % arancio - esempi
\definecolor{remcol}{HTML}{6B7280}   % grigio  - osservazioni

% ---------- Listings (riusa la palette nuova) ----------
\usepackage{listings}
\lstdefinestyle{mystyle}{
    backgroundcolor=\color{remcol!8},
    commentstyle=\color{defcol},
    keywordstyle=\color{thmcol},
    numberstyle=\tiny\color{remcol},
    stringstyle=\color{lemcol},
    basicstyle=\ttfamily\footnotesize,
    breakatwhitespace=false, breaklines=true, captionpos=b,
    keepspaces=true, numbers=left, numbersep=5pt,
    showspaces=false, showstringspaces=false, showtabs=false, tabsize=2
}
\lstdefinestyle{terminal}{
    backgroundcolor=\color{black},
    basicstyle=\ttfamily\color{green},
    frame=single, breaklines=true, showstringspaces=false
}
\lstset{style=mystyle}

% ---------- Teoremi, definizioni, esempi (tcolorbox) ----------
\usepackage[most]{tcolorbox}
\tcbuselibrary{theorems}

\tcbset{
  thmstyle/.style={
    enhanced, breakable, sharp corners,
    boxrule=0pt, frame hidden,
    borderline west={3pt}{0pt}{#1},
    colback=#1!6,
    colbacktitle=#1!14, coltitle=#1!60!black,
    fonttitle=\bfseries,
    titlerule=0pt,
    separator sign none,
    description delimiters parenthesis,
    before skip=12pt, after skip=12pt,
    left=8pt, right=8pt, top=4pt, bottom=4pt,
    overlay unbroken and first={\tikzset{external/export=false}},
    overlay middle and last={\tikzset{external/export=false}},
  }
}

\newtcbtheorem[number within=section]{theorem}{Teorema}{thmstyle=thmcol}{thm}
\newtcbtheorem[use counter from=theorem]{lemma}{Lemma}{thmstyle=lemcol}{lem}
\newtcbtheorem[use counter from=theorem]{definition}{Definizione}{thmstyle=defcol}{def}
\newtcbtheorem[use counter from=theorem]{example}{Esempio}{thmstyle=excol}{ex}
\newtcolorbox{remark}[1][]{thmstyle=remcol, title={Osservazione}, #1}

% ---------- Esercizi (mdframed) ----------
\usepackage{mdframed}
\theoremstyle{definition}
\newmdtheoremenv[
    linecolor=gray, leftmargin=0, rightmargin=0,
    innertopmargin=8pt, innerbottommargin=8pt
]{exercise}{Esercizio}[section]

% ---------- Figure svg ----------
\newcommand{\incfig}[1]{%
    \def\svgwidth{\columnwidth}
    \import{./figures/}{#1.pdf_tex}
}

% ---------- Hyperref (sempre per ultimo, prima di externalize) ----------
\usepackage{hyperref}
\hypersetup{colorlinks=false}

% ---------- Externalize ----------
\usetikzlibrary{external}
\tikzexternalize[prefix=figures/tikz/, optimize=false]
\tikzexternaldisable



\begin{document}
\begin{titlepage}
	\begin{center}
		\vspace*{1cm}

		\Huge
		\textbf{Natural Language Processing}

		\vspace{0.5cm}
		\LARGE
		UniVR - Dipartimento di Informatica

		\vspace{1.5cm}

		\textbf{Tommaso Toffali}

		\vfill


		\vspace{0.8cm}


		1° Semestre 2026/2027
    \end{center}
\end{titlepage}

\tableofcontents
\pagebreak

\part{String Similarity and Edit Distance}
\section{Minimum Edit Distance}
We are going to deal with this main driving point: the definition of Minimum Edit Distance.
\vspace{1em}
\noindent
\newline
\textbf{The question: are this 2 texts the same?}
\vspace{1em}
\noindent
\newline
When is the case when 2 texts are the same? Of course when every single character is the same.
What if we would like to understand if 2 texts are pretty close (not the same)?
\vspace{1em}
\noindent
\newline
Single characters in the text could be different in position.
\subsection{How similar are two strings?}
The fundamental question in edit distance is: how can we measure the similarity between two 
strings? This problem appears in many different applications across various domains of computer 
science and computational linguistics.
\subsubsection{Spell Correction}
In spell correction systems, we need to find the closest valid word to a misspelled input. For 
example, if the user typed "graffe", we need to determine which word in our dictionary is most similar.
\vspace{1em}
\noindent
\newline
Possible candidates:
\begin{itemize}
  \item graf 
  \item graft 
  \item grail 
  \item giraffe
\end{itemize}
By computing the edit distance between "graffe" and each candidate, we can identify that "giraffe" is likely the intended 
word (requiring only one deletion).
\subsubsection{Computational Biology}
In bioinformatics, sequence alignment is crucial for comparing DNA, RNA, or protein sequences.
By aligning sequences, we can identify regions of similarity that may indicate functional, structural, or evolutionary relationships.
\vspace{1em}
\noindent
\newline
Example: Align two sequences of nucleotides:
\vspace{1em}
\noindent
\newline
\texttt{AGGCTATCACCTGACCTCCAGGCCGATGCCC}
\vspace{1em}
\noindent
\newline
\texttt{TAGCTATCACGACCGCGGTCGATTTGCCCGAC}
\vspace{1em}
\noindent
\newline
Resulting alignment (dashes represent insertions/deletions):
\vspace{1em}
\noindent
\newline
\texttt{-AGGCTATCACCTGACCTCCAGGCCGA--TGCCC---}
\newline
\texttt{TAG-CTATCAC--GACCGC--GGTCGATTTGCCCGAC}
\vspace{1em}
\noindent
\newline
The alignment reveals matching regions (shown in the same positions) and differences between the sequences.
\subsubsection{Other Applications}
String similarity and edit distance algorithms are fundamental tools in many NLP tasks:
\begin{itemize}
  \item \textbf{Machine Translation}: Comparing source and target language phrases, finding similar translations in translation memories.
  \item \textbf{Information Extraction}: Matching entity names with variations (e.g., "IBM" vs "I.B.M." vs "International Business Machines")
  \item \textbf{Speech Recognition}: Correcting recognition errors by finding the closest valid word or phrase to the acoustic model output
\end{itemize}
\subsection{Edit Distance}
The \textbf{minimum edit distance} between two strings is defined as the minimum number of editing operations needed to transform 
one string into the other
\subsubsection{Edit Operations}
\begin{definition}{Three Basic Edit Operations}{Three Basic Edit Operation}
There are three basic editing operations:
\begin{itemize}
  \item \textbf{Insertion}: Add a character at any position in the string 
  \begin{itemize}
    \item Example: cat \( \rightarrow \) cart (insert 'r')
  \end{itemize}
  \item \textbf{Deletion}: Remove a character from any position in the string 
  \begin{itemize}
    \item Example: cart \( \rightarrow \) cat (delete 'r')
  \end{itemize}
  \item \textbf{Substitution}: Replace one character with another 
  \begin{itemize}
    \item Example: cat \( \rightarrow \) bat (substitute 'c' with 'b')
  \end{itemize}
\end{itemize}
\subsubsection{Key Concept}
The edit distance measures the \textit{minimum} number of these operations required to transform one 
string into another. This metric provides a quantitative measure of string similarity: the smaller 
the edit distance, the more similar the strings are.
\vspace{1em}
\noindent
\newline
\textbf{Important note}: Each operation has a cost (typically 1), and we seek the sequence of operations
that minimizes the total cost. Different variants of edit distance may assign different costs 
to different operations (e.g., Levenshtein distance uses unfirom costs, while other variants may weight substitutions differently)
\end{definition}
\subsubsection{Example: Alignment of Two Strings}
To better understand edit distance, let's examine how two strings can be aligned to show their differences.
Consider the string "INTENTION" and "EXECUTION":
\begin{center}
  \setlength{\tabcolsep}{7pt}   % spazio tra le colonne
  \begin{tabular}{*{10}{c}}
    I & N & T & E & $*$ & N & T & I & O & N \\
    $|$ & $|$ & $|$ & $|$ & $|$ & $|$ & $|$ & $|$ & $|$ & $|$ \\
    $*$ & E & X & E & C & U & T & I & O & N
  \end{tabular}
\end{center}
In this alignment:
\begin{itemize}
  \item The asterisk (*) represents a gap, indicating an insertion or deletion operation 
  \item Vertical bars (|) connect corresponding positions between the two strings 
  \item Characters that match are aligned vertically (E, T, I, O, N)
  \item Characters that differ indicate a substitution operations
\end{itemize}
\textbf{Operations needed to transform "INTENTION" to "EXECUTION"}
\begin{enumerate}
  \item Delete 'I' at position 1 
  \item Substitute 'N' with 'E' at position 2 
  \item Substitute 'T' with 'X' at position 3
  \item Keep 'E' (match)
  \item Insert 'C' at position 5
  \item Substitute 'N' with 'U' at position 6 
  \item Keep 'T' (match)
  \item Keep 'I' (match)
  \item Keep 'O' (match)
  \item Keep 'N' (match)
\end{enumerate}
This gives us a total edit distance of \textbf{5 operations} (1 deletion + 3 substitution + 1 insertion)
\subsubsection{Cost variants: Standard vs Levenshtein}
The total distance depends on how we assign costs to each operation:
\begin{definition}{Standard Edit Distance (uniform cost)}{Standard Edit Distance (uniform cost)}
\begin{itemize}
  \item Each operation (insertion, deletion, substitution) costs 1 
  \item Total distance for INTENTION \( \rightarrow \) EXECUTION: \( \mathbf{5} \)   
  \item Calculation: 1 (deletion) + 3 \( \times \) 1 (substitutions) + 1 (insertion) = 5
\end{itemize}
\end{definition}
\begin{theorem}{Levenshtein Distance (weighted substitution)}{Levenshtein Distance (weighted substitution)}
\begin{itemize}
  \item Insertion costs 1 
  \item Deletion costs 1 
  \item Substitution costs 2 (considered as a deletion + insertion, hence a "double error")
  \item Total distance for INTENTION \( \rightarrow \) EXECUTION: \( \mathbf{8} \)    
  \item Calculation: 1 (deletion) + 3 \( \times \) 2 (substitution) + 1 (insertion) = 8
\end{itemize}
\end{theorem}
\textbf{Why the difference?} In the Levenshtein variant, a substitution is viewed as conceptually 
equivalent to deleting a character and then inserting a different one, thus costing twice as much.
This distinction is important when choosing which distance metric to use for a particular relationships.
\subsection{Alignment in Computational Biology}
In computation biology, sequence alignment is a fundamental technique for comparing DNA, RNA, or protein sequences.
The goal is to identify regions of similarity and understand evolutionary relationships.
\subsubsection{The alignment problem}
\textbf{Given a sequence of bases:}
\vspace{1em}
\noindent 
\newline 
\texttt{AGGCTATCACCTGACCTCCAGGCCGATGCCC}
\vspace{1em}
\noindent 
\newline
\texttt{TAGCTATCACGACCGCGGTCGATTTGCCCCGAC}
\vspace{1em}
\noindent 
\newline
\textbf{An alignment}:
\vspace{1em}
\noindent 
\newline
\texttt{-AGGCTATCACCTGACCTCCAGGCCGA---TGCCC---}
\newline 
\texttt{TAG-CTATCAC--GACCGC--GGTCGATTTGCCCCGAC}
\subsubsection{Alignment Objective}
\textbf{Given two sequences, align each letter to a letter or gap}
\vspace{1em}
\noindent 
\newline
The alignment process involves:
\begin{itemize}
  \item \textbf{Matching}: Aligning identical bases (e.g., A with A, G with G)
  \item \textbf{Mismatches}: Aligning different bases (substitutions)
  \item \textbf{Gaps}: Represented by dashes (-), indicating insertions or deletions (indels)
\end{itemize}
\textbf{Key considerations}:
\begin{itemize}
  \item The alignment should maximize the number of matches 
  \item Minimize the number of mismatches and gaps 
  \item Gaps are penalized because schemes can be used: match scores, mismatch penalties, and gap penalites
\end{itemize}
This alignment problem is directly related to edit distance: finding the optimal alignment is equivalent 
to finding the minimum edit distance between the two sequences.
\subsection{Other uses of Edit distance in NLP}
Edit distance is versatile tool used across many NLP applications beyond spell correction and 
sequence alignment.
\subsubsection{Evaluating Machine Translation and Speech Recognition}
Edit distance can be used to evaluate the quality of machine translation and speech recognition 
systems by comparing the system output with a reference (correct) translation or transcription.
\vspace{1em}
\noindent 
\newline
\textbf{Example}:
\vspace{1em}
\noindent 
\newline
\textbf{R} (Reference): Spokesman confirms senior government adviser was appointed 
\vspace{1em}
\noindent 
\newline
\textbf{H} (Hypothesis): Spokesman said the senior adviser was appointed 
\[
  \text{S} \quad \text{I} \quad \text{D} \quad \text{I}
\]
Where:
\begin{itemize}
  \item \textbf{S} = Substitution ("confirms" $\rightarrow$ "said")
  \item \textbf{I} = Insertion ("the" inserted)
  \item \textbf{D} = Deletion ("government" deleted)
  \item \textbf{I} = Insertion (extra word)
\end{itemize}
The edit distance provides a quantitative measure of how different the hypothesis is from the 
reference, which is crucial for evaluating system performance.
\subsubsection{Named Entity Extraction and Entity Coreference}
Edit distance helps identify when different text strings refer to the same entity, even when they 
are written differently.
\vspace{1em}
\noindent 
\newline
\textbf{Examples}:
\begin{itemize}
  \item
\end{itemize}
\end{document}